\section{Conclusion}
\label{sec:conclusion}

A weather forecast for renewable energy is good only relative to a decision. Fixing the technology, the department and the horizon determines which variables matter, at what resolution and by when, in which uncertainty format, and with which scoring rule. A PV trader needs quarter-hourly irradiance quantiles before noon; a crane supervisor needs gusts at hub height for the next morning; a bank needs the P90 of twenty years. Four results govern the design of every such system: the power curve amplifies wind errors by a factor of at least three at partial load; errors shared across sites cannot be diversified away; averaging two models helps only when they are similarly good and fail differently; and under asymmetric imbalance costs the right bid is a quantile, not the expected value.

The case study confirms all four on public data and adds two practical findings. A machine-learning model, ECMWF's AIFS, now matches its physics-based counterpart for day-ahead station wind and beats it for temperature, while NOAA GFS trails both. And the largest gains in the study came not from choosing a model but from what follows it: averaging three models and removing a running bias extended the useful daytime wind forecast from about five to about eight and a half days and cut the day-ahead error by about a sixth. The data, code and tables needed to repeat or extend the analysis are published with this paper; the obvious next steps are hub-height wind, measured irradiance and plant output, and a winter season.
